\begin{frame}{RLHF Training: Proximal Policy Optimization (PPO)}
  \begin{itemize}
    \item[\textcolor{red}{$\blacktriangleright$}] \textbf{Objective:} Maximize the reward given by the reward model (RM).
    \item[\textcolor{red}{$\blacktriangleright$}] \textbf{Algorithm:} Proximal Policy Optimization (PPO) is used to update the language model.
    \item[\textcolor{red}{$\blacktriangleright$}] \textbf{PPO Loss:}
      \[
        L_{\mathrm{PPO}} = \mathbb{E}\left[
            \min\Big(r_t(\theta)\hat{A}_t,\,\,\,\, \text{clip}\big(r_t(\theta), 1 - \epsilon, 1 + \epsilon\big)\hat{A}_t\Big)
          \right]
      \]
      where $r_t(\theta) = \frac{\pi_\theta(a_t|s_t)}{\pi_\text{old}(a_t|s_t)}$ is the probability ratio, and $\hat{A}_t$ is the advantage estimate.
    \item[\textcolor{red}{$\blacktriangleright$}] \textbf{KL Penalty:} A KL-divergence penalty is added to keep the policy close to the base (SFT) model:
      \[
        L_{\mathrm{KL}} = \beta\, \mathrm{KL}\Big(\pi_\theta \;\|\; \pi_{\text{SFT}}\Big)
      \]
  \end{itemize}
\end{frame}


\begin{frame}{RLHF Training: Proximal Policy Optimization (PPO)}
  \centering
  \includegraphics[width=0.60\paperwidth,height=0.60\paperheight,keepaspectratio]{assets/Lecture-8_Fine_Tuning_RLHF/slide_022.png}
\end{frame}

\begin{frame}{PPO --- Pros and Cons}
  \begin{itemize}
    \item[\textcolor{red}{\checkmark}] Stable optimization
    \item[\textcolor{red}{\checkmark}] Encourages exploration
    \item[\textcolor{gray}{\texttimes}] Expensive (needs many rollouts)
    \item[\textcolor{gray}{\texttimes}] Sensitive to reward model errors
  \end{itemize}
\end{frame}
